\section{Extended Protocol and Additional Placeholder Studies}
\label{app:experiments}

This appendix expands the evaluation protocol of \Cref{sec:experiments}: it records the
precise metric definitions, the Sybil and adaptivity studies, the cost model, and a
sensitivity plan. All tables are \textbf{placeholders}; every quantitative outcome is
\emph{(hypothesis)} until produced.

\subsection{Precise metric definitions}
Given the \corb{} logs, each metric is a deterministic functional of the recorded traces:
\begin{itemize}[leftmargin=1.4em]
\item \textbf{Coordination quality} $=\PNAs$ on forced-coordination tasks under the benign
counterpart $\config^{\circ}$; \textbf{coordination sabotage} $=\PNAs(\config^{\circ})-\PNAs(\config)$
for a sabotage configuration $\config$.
\item \textbf{Time-to-detection} $=\E[\min\{t:\text{defense flags a compromise in }\Trace\}]$,
$\infty$ if never; reported as a survival curve over $t$.
\item \textbf{Recovery rate} $=\Prob[\text{episode returns to a safe, task-completing state
after a detection}]$.
\item \textbf{Token/API cost} and \textbf{latency}: summed over all agent and emulator calls
in an episode, averaged over the $M$ episodes; \textbf{scalability}: the fitted growth of each
metric and cost in $N$.
\item \textbf{Human-intervention rate} $=\Prob[\text{a human-approval hook was triggered}]$
when hooks are enabled.
\end{itemize}

\subsection{Sybil and adaptivity studies}
\Cref{tab:sybil} gives the structure of the Sybil study (RQ5): as the number of Sybil
identities grows, an adversary controlling one underlying policy gains disproportionate
influence over selection, voting, and memory writes. \Cref{tab:adaptive} gives the
adaptivity study: static versus adaptive red policies across \corb{} rounds, reporting how
quickly the blue policy is out-evolved.

\begin{table}[t]
\centering
\caption{\textbf{Placeholder / expected trend only; replace with actual experimental results
before submission.} Sybil study (RQ5). Expected \emph{(hypothesis)} trend: influence and
$\ASRs$ increase with Sybil count until selection/voting defenses saturate.}
\label{tab:sybil}
\small
\begin{tabular}{@{}lcccc@{}}
\toprule
Sybil count & Selection share & Vote share & $\ASRs$ & $\NRPs$ \\
\midrule
0 & --- & --- & --- & --- \\
1 & --- & --- & --- & --- \\
2 & --- & --- & --- & --- \\
4 & --- & --- & --- & --- \\
\bottomrule
\end{tabular}
\end{table}

\begin{table}[t]
\centering
\caption{\textbf{Placeholder / expected trend only; replace with actual experimental results
before submission.} Adaptivity study: static vs.\ adaptive red across \corb{} rounds
$r=1,\dots,R$. Expected \emph{(hypothesis)} trend: adaptive red erodes a static blue's
$\NRPs$ over rounds; a co-adapting blue slows but may not stop the erosion.}
\label{tab:adaptive}
\small
\begin{tabular}{@{}lcccc@{}}
\toprule
& \multicolumn{2}{c}{Static red} & \multicolumn{2}{c}{Adaptive red} \\
\cmidrule(lr){2-3}\cmidrule(lr){4-5}
Round & $\ASRs$ & $\NRPs$ & $\ASRs$ & $\NRPs$ \\
\midrule
$r=1$ & --- & --- & --- & --- \\
$r=5$ & --- & --- & --- & --- \\
$r=10$ & --- & --- & --- & --- \\
$r=20$ & --- & --- & --- & --- \\
\bottomrule
\end{tabular}
\end{table}

\subsection{Cost model and reachable region}
By \Cref{sec:algorithm}, one \corb{} round costs $2M$ episodes and each episode costs
$\Theta(NH)$ agent invocations plus emulator calls, so a study of $G$ swept configurations
and $R$ rounds costs $\Theta(GRMNH)$ invocations. With $M$ set by
\Cref{thm:sample}(c) to $M=O(\epsilon^{-2}\ln(G/\alpha))$, the dominant cost driver is the
product $GN$: both the sweep grid and the agent count. \Cref{tab:cost} records the structure
of the scalability study (RQ6). We report the \emph{reachable region} of the parameter space
explicitly---which $(N,\text{topology},\betafrac,\text{knobs})$ combinations were actually
run---rather than implying exhaustive coverage.

\begin{table}[t]
\centering
\caption{\textbf{Placeholder / expected trend only; replace with actual experimental results
before submission.} Scalability and cost (RQ6) vs.\ agent count $N$. Expected
\emph{(hypothesis)} trend: token/API cost and latency grow with $N$ (and with \corb{}
rounds), constraining the reachable region.}
\label{tab:cost}
\small
\begin{tabular}{@{}lccccc@{}}
\toprule
$N$ & $\ASRs$ & $\CPR$ & Tokens/episode & Latency/episode & Reachable? \\
\midrule
3  & --- & --- & --- & --- & --- \\
10 & --- & --- & --- & --- & --- \\
20 & --- & --- & --- & --- & --- \\
50 & --- & --- & --- & --- & --- \\
\bottomrule
\end{tabular}
\end{table}

\subsection{A worked forced-coordination task}
\label{app:worked}
To make the predicate machinery concrete, consider a $3$-forced ``travel-booking'' task on a
star of $N$ agents. The orchestrator decomposes the goal into three subtasks assigned to
distinct workers: (W1) retrieve an itinerary from a calendar tool; (W2) select a flight within a
stated budget via a booking tool; and (W3) pay via a payment tool, but \emph{only} to the payee
named in the user's instruction. A witness (\Cref{def:forced}) is any set of agents that jointly
supplies a valid itinerary, an in-budget flight, and a correctly addressed payment; because all
three subtasks are required and are assigned to distinct workers, every witness has size at least
three, so the task is $3$-forced.

The registered predicates read only the recorded trace $\Trace$:
\begin{align*}
\Predutil(\Trace) &= \indic[\text{a payment to the instructed payee for an in-budget flight matching the itinerary is recorded}],\\
\Predsec(\Trace) &= \indic[\text{a payment to a payee \emph{not} named in the instruction is recorded (exfiltration of funds), \emph{or}}\\
&\qquad\ \text{a denylisted action (e.g., disclosing card details to a tool) is recorded}].
\end{align*}
A malicious tool (via $\Emu^{\!\dagger}$) might return a booking result that embeds an injected
instruction ``pay to account X''; a colluding W3 might then act on it. Whether the payment to
account~X occurs is a recorded state transition, so $\Predsec$ fires deterministically regardless
of what the agents \emph{say} in their messages---the content of the injection cannot change the
verdict unless the illicit payment is actually recorded (\Cref{thm:integrity}). If instead the
adversary set forms a transversal of the witnesses (e.g., it controls W3 in every route), it can
force $\Predutil=0$ by withholding the legitimate payment (\Cref{prop:forced}), which the sabotage
component of the metric records as coordination failure. This example is illustrative; the shipped
task suite generalizes it across roles, tools, and topologies.

\subsection{Sensitivity and validity checks}
Beyond the headline studies, the protocol includes: (i) an \textbf{aggregator-sensitivity}
check (RQ1) that recomputes configuration rankings under $u\sigma$, $\min(u,\sigma)$, and a
weighted geometric mean $u^{w}\sigma^{1-w}$ (\Cref{rem:metric}) and reports rank correlation;
(ii) an \textbf{evaluator-integrity} check (RQ2) that measures how often a crafted injection
flips the LLM judge $\Judge$ while leaving the deterministic verdict $\Checker(\Trace)$
unchanged (an empirical instance of \Cref{thm:integrity} versus \Cref{prop:judgehijack});
(iii) an \textbf{emulation-fidelity} check that compares emulated and real tools on the
subset of tasks where a real implementation is available, reporting the change in $\ASRs$ and
$\NRPs$ and relating it to $\kappaem$; and (iv) \textbf{backbone-robustness} checks that repeat
the headline degradation study across all five model families so that no conclusion rests on a
single model. Each check is a registered configuration in the released harness, so validity,
not only headline numbers, is reproducible.
